\documentclass[10pt, letterpaper]{article}

% Packages:
\usepackage[
    ignoreheadfoot, % set margins without considering header and footer
    top=0.3 cm, % seperation between body and page edge from the top
    bottom=0.2 cm, % seperation between body and page edge from the bottom
    left=0.2 cm, % seperation between body and page edge from the left
    right=0.2 cm, % seperation between body and page edge from the right
    footskip=0 cm, % seperation between body and footer
    % showframe % for debugging 
]{geometry} % for adjusting page geometry
\usepackage{titlesec} % for customizing section titles
\usepackage{tabularx} % for making tables with fixed width columns
\usepackage{array} % tabularx requires this
\usepackage[dvipsnames]{xcolor} % for coloring text
\definecolor{primaryColor}{RGB}{0, 79, 144} % define primary color
\usepackage{enumitem} % for customizing lists
\usepackage{fontawesome5} % for using icons
\usepackage{amsmath} % for math
\usepackage[
    pdftitle={Nikhil Pise's CV},
    pdfauthor={Nikhil Pise},
    pdfcreator={},
    colorlinks=true,
    urlcolor=primaryColor
]{hyperref} % for links, metadata and bookmarks
\usepackage[pscoord]{eso-pic} % for floating text on the page
\usepackage{calc} % for calculating lengths
\usepackage{bookmark} % for bookmarks
\usepackage{lastpage} % for getting the total number of pages
\usepackage{changepage} % for one column entries (adjustwidth environment)
\usepackage{paracol} % for two and three column entries
\usepackage{ifthen} % for conditional statements
\usepackage{needspace} % for avoiding page brake right after the section title
\usepackage{iftex} % check if engine is pdflatex, xetex or luatex

% Ensure that generate pdf is machine readable/ATS parsable:
\ifPDFTeX
    \input{glyphtounicode}
    \pdfgentounicode=1
    % \usepackage[T1]{fontenc} % this breaks sb2nov
    \usepackage[utf8]{inputenc}
    \usepackage{lmodern}
\fi



% Some settings:
\AtBeginEnvironment{adjustwidth}{\partopsep0pt} % remove space before adjustwidth environment
\pagestyle{empty} % no header or footer
\setcounter{secnumdepth}{0} % no section numbering
\setlength{\parindent}{0pt} % no indentation
\setlength{\topskip}{0pt} % no top skip
\setlength{\columnsep}{0cm} % set column seperation
\makeatletter


\titleformat{\section}{\needspace{4\baselineskip}\bfseries\large}{}{0pt}{}[\vspace{1pt}\titlerule]

\titlespacing{\section}{
    % left space:
    0.1 cm
}{
    % top space:
    0.2 cm
}{
    % bottom space:
    0.2 cm
} % section title spacing

\renewcommand\labelitemi{$\circ$} % custom bullet points
\newenvironment{highlights}{
    \begin{itemize}[
        topsep=0.10 cm,
        parsep=0.10 cm,
        partopsep=0pt,
        itemsep=0pt,
        leftmargin=0.4 cm + 10pt
    ]
}{
    \end{itemize}
} % new environment for highlights

\newenvironment{highlightsforbulletentries}{
    \begin{itemize}[
        topsep=0.10 cm,
        parsep=0.10 cm,
        partopsep=0pt,
        itemsep=0pt,
        leftmargin=10pt
    ]
}{
    \end{itemize}
} % new environment for highlights for bullet entries


\newenvironment{onecolentry}{
    \begin{adjustwidth}{
        0.2 cm + 0.00001 cm
    }{
        0.2 cm + 0.00001 cm
    }
}{
    \end{adjustwidth}
} % new environment for one column entries

\newenvironment{twocolentry}[2][]{
    \onecolentry
    \def\secondColumn{#2}
    \setcolumnwidth{\fill, 4.5 cm}
    \begin{paracol}{2}
}{
    \switchcolumn \raggedleft \secondColumn
    \end{paracol}
    \endonecolentry
} % new environment for two column entries

\newenvironment{header}{
    \setlength{\topsep}{0pt}\par\kern\topsep\centering\linespread{1.5}
}{
    \par\kern\topsep
} % new environment for the header

\newcommand{\placelastupdatedtext}{% \placetextbox{<horizontal pos>}{<vertical pos>}{<stuff>}
  \AddToShipoutPictureFG*{% Add <stuff> to current page foreground
    \put(
        \LenToUnit{\paperwidth-2 cm-0.2 cm+0.05cm},
        \LenToUnit{\paperheight-1.0 cm}
    ){\vtop{{\null}\makebox[0pt][c]{
        \small\color{gray}\textit{Last updated in January 2025}\hspace{\widthof{Last updated in January 2025}}
    }}}%
  }%
}%

% save the original href command in a new command:
\let\hrefWithoutArrow\href

% new command for external links:
\renewcommand{\href}[2]{\hrefWithoutArrow{#1}{\ifthenelse{\equal{#2}{}}{ }{#2 }\raisebox{.15ex}{\footnotesize \faExternalLink*}}}


\begin{document}
    \newcommand{\AND}{\unskip
        \cleaders\copy\ANDbox\hskip\wd\ANDbox
        \ignorespaces
    }
    \newsavebox\ANDbox
    \sbox\ANDbox{}

    \begin{header}
        \textbf{\fontsize{24 pt}{24 pt}\selectfont Nikhil Pravin Pise}

        \vspace{0 cm}
        \normalsize
        
        \mbox{\hrefWithoutArrow{tel:+91-9867477169}{\color{black}{\footnotesize\faPhone*}\hspace*{0.13cm}\href{tel:+91-9867477169/}{Telephone}}}%
        \AND%        
        \kern 0.25 cm%
        \mbox{\hrefWithoutArrow{mailto:nikhilpise2006@gmail.com}{\color{black}{\footnotesize\faEnvelope[regular]}\hspace*{0.13cm}\href{mailto:nikhilpise2006@gmail.com}{Gmail} }}%
        \AND%        
        \kern 0.25 cm%
        \mbox{\hrefWithoutArrow{www.linkedin.com/in/nikhil-pravin-pise}{\color{black}{\footnotesize\faLinkedinIn}\hspace*{0.13cm}\href{https://www.linkedin.com/in/nikhil-pravin-pise}{LinkedIn} }}%
        \AND%        
        \kern 0.25 cm%
        \mbox{\hrefWithoutArrow{https://github.com/N1KH1LT0X1N}{\color{black}{\footnotesize\faGithub}\hspace*{0.13cm}\href{https://github.com/N1KH1LT0X1N}{GitHub} }}%
        \AND%        
        \kern 0.25 cm%
        \mbox{\hrefWithoutArrow{https://www.nikhilpise.me/}{\color{black}{\footnotesize\faLink}\hspace*{0.13cm}\href{https://www.nikhilpise.me/}{Portfolio} }}%
        
    \end{header}

    \vspace{0 cm}

    \section{Education} 
    
        \begin{twocolentry}{
            \textit{Sept 2024 – Present}            
        
            \textit{9.263 CGPA}}
            \textbf{SVKM's Dwarkadas Jivanlal Sanghvi College of Engineering}    
            
            \textit{Computer Science and Engineering : Data Science}
        \end{twocolentry}
        
        %\begin{twocolentry}{
        %    \textit{Aug 2022 – Aug 2024}    
        %
        %    \textit{88.17 \%}}
        %    \textbf{Pace Junior Science College of Engineering}      
        %    
        %    \textit{Computer Science, Physics, Chemistry, Mathematics, English}
        %\end{twocolentry}

        %\begin{twocolentry}{
        %    \textit{Jul 2013 – Jul 2022}    
        %    
        %    \textit{92.50 \%}}
        %    \textbf{Thakur International School}  
        %    
        %    \textit{Computer Applications, Physics, Chemistry, Biology, Mathematics, %English, Hindi}
        %\end{twocolentry}

                     
    \section{Achievements}

        \begin{onecolentry}
            \begin{highlightsforbulletentries}
                \item \text{\textbf{\href{https://drive.google.com/file/d/1kr2mnXLC7Cfrm538Si7AMdGfPy8dEBxB/view?usp=drive_link}{Technograd 4.0 – 2nd Runner-Up \faTrophy} – DJS NSDC}, Nov 2025.}
                \item \text{\textbf{\href{https://drive.google.com/file/d/1AHYOCVgGzDN7_ScU8IZsH-XjbTjfEaOH/view?usp=drive_link}{MCP's 1st Birthday – 4th Place \faTrophy} – Gradio x Anthropic}, Dec 2025}
                \item \text{\textbf{\href{https://drive.google.com/file/d/1Ru4X2Lx2TTQcC44BP1JUSrxGD6nRhP14/view?usp=drive_link}{Megathon 2025 – Top 10 Finalist \faMedal} – IIIT Hyderabad}  $|$}                \text{\textbf{\href{https://drive.google.com/file/d/1wB0g95G_iXeF6WbTxuaXVo36Wys5MUAx/view?usp=drive_link}{Redact 2025 – Top 13 Finalist \faMedal} –SPIT Mumbai}}
                \item \text{\textbf{\href{https://drive.google.com/file/d/1H-NXnbKrQZ8eEwLXcwaeaO87uFzI8Pgb/view?usp=drive_link}{3 Preliminary Astroid Discoveries} – IASC x Nasa}, Dec 2025.}
            \end{highlightsforbulletentries}
        \end{onecolentry}     
       

    \section{Experience}
    
        \begin{twocolentry}{
        \textit{Mumbai}    
            
        \textit{Sep 2025 - Present}}
            \textbf{ML Co-Committee Member}
            
            \textit{DJS S4DS}
        \end{twocolentry}

        \vspace{0.10 cm}
        \begin{onecolentry}
            \begin{highlights}
                \item \textbf{FIFA Player Analytics (\href{https://github.com/N1KH1LT0X1N/Plotting}{GitHub}):} Player salaries lacked data-driven benchmarks, so I analyzed and feature-engineered \textbf{7 FIFA seasons (2017--2023)} to build explainable wage-prediction models and interactive dashboards for smarter player valuation.
                \item \textbf{Extract 4.0 Hackathon:} The competition needed a truly differentiating challenge, so I authored \textbf{1 hard-difficulty} problem with edge case datasets and an objective scoring rubric to separate top problem-solvers from the rest.
            \end{highlights}
        \end{onecolentry}


    

    
    \section{Projects}     
                 
        %\begin{twocolentry}{                     
        %    \textit{2025}}
        %    \textbf{Traffic | Python, TensorFlow, OpenCV, CNN}
        %\end{twocolentry}
        %
        %\vspace{0.0 cm}
        %
        %\begin{onecolentry}
        %    \begin{highlights}
        %        \item Built a convolutional neural network to classify 43 different %German traffic signs with high accuracy using TensorFlow and OpenCV.
        %        \item Achieved high accuracy in CV classification tasks, demonstrating %practical AI applications for autonomous vehicle safety systems.
        %    \end{highlights}
        %\end{onecolentry}        

        %\vspace{0.1 cm}

        %\begin{twocolentry}{                     
        %    \textit{2025}}
        %    \textbf{Questions | Python, NLP, TF-IDF}
        %\end{twocolentry}

        %\vspace{0.0 cm}
        
        %\begin{onecolentry}
        %    \begin{highlights}
        %        \item Built a qna system that retrieves documents and extracts answers %using TF-IDF ranking and query term density analysis.
        %        \item Demonstrated information retrieval and natural language processing %capabilities for automated document analysis.
        %    \end{highlights}
        %\end{onecolentry}

        %\vspace{0.1 cm}

        \begin{twocolentry}{
\textit{Open-source · Production}}
\textbf{Medium Agent / Project Aether (\href{https://huggingface.co/spaces/T0X1N/Medium-MCP}{Hugging Face})}
\end{twocolentry}

\vspace{0.1 cm}

\begin{onecolentry}
\begin{highlights}
    \item Medium content was hard to search and consume at scale, so I built \textbf{Project Aether}, a Gradio-based research and Sonic-Mode (neural podcast) interface that turns Medium into a \textbf{queryable and listenable knowledge base}.
    \item LLMs couldn’t reason over Medium autonomously, so I built a production \textbf{MCP v3.0 server} exposing \textbf{12 tools and 4 resources}, enabling \textbf{Claude Desktop} and other agents to search, scrape, and synthesize articles end-to-end.
    \item \textbf{Tech:} \textbf{MCP}, \textbf{Groq LLaMA}, \textbf{Gemini}, \textbf{ElevenLabs}, \textbf{OpenAI}, \textbf{Gradio}, \textbf{Playwright}, \textbf{Docker}, \textbf{Python}
\end{highlights}
\end{onecolentry}

\vspace{0.1 cm}

\begin{twocolentry}{
\textit{Open-source}}
\textbf{Pramana.ai — Biomedical Research Platform (\href{https://github.com/N1KH1LT0X1N/Pramana.AI}{GitHub})}
\end{twocolentry}

\vspace{0.1 cm}

\begin{onecolentry}
\begin{highlights}
    \item Biomedical research was slow and fragmented, so I built a \textbf{LangGraph-based multi-agent system} that runs \textbf{5 specialized AI agents} for drug discovery, evidence review, and report generation.
    \item Multi-agent reasoning was opaque, so I added \textbf{real-time execution graphs} and a visual workflow canvas, making every step of model reasoning inspectable and debuggable.
    \item \textbf{Tech:} \textbf{LangGraph}, \textbf{LangChain}, \textbf{Groq (LLaMA 3.3 70B)}, \textbf{FastAPI}, \textbf{React 19}, \textbf{Redis}, \textbf{PostgreSQL}, \textbf{Docker}, \textbf{MCP}
\end{highlights}
\end{onecolentry}

\vspace{0.1 cm}

\begin{twocolentry}{
\textit{Open-source · Live}}
\textbf{MediGuard AI — Intelligent Triage Platform (\href{https://github.com/N1KH1LT0X1N/MediGuard}{GitHub})}
\end{twocolentry}

\vspace{0.1 cm}

\begin{onecolentry}
\begin{highlights}
    \item Blood test interpretation was manual and error-prone, so I built an \textbf{AI triage engine} using \textbf{24 biomarkers} to predict \textbf{6+ diseases} with \textbf{99.12\% accuracy} and explainable probability scores.
    \item Patients lacked real-time access to clinical insight, so I deployed a \textbf{WhatsApp medical assistant} for instant blood-test analysis, backed by \textbf{PubMed, WHO, and AHA} with anonymized audit logs (\href{https://mediguard-chatbot-frontend.vercel.app/}{Live}).
    \item \textbf{Tech:} \textbf{XGBoost}, \textbf{FastAPI}, \textbf{React}, \textbf{Flask}, \textbf{Gemini}, \textbf{Twilio}, \textbf{PostgreSQL}, \textbf{Docker}, \textbf{Python}
\end{highlights}
\end{onecolentry}

\vspace{0.1 cm}

\begin{twocolentry}{
\textit{Open-source · Production}}
\textbf{Spotify MCP Server (\href{https://github.com/N1KH1LT0X1N/SpotifyMCP}{GitHub})}
\end{twocolentry}

\vspace{0.1 cm}

\begin{onecolentry}
\begin{highlights}
    \item Controlling Spotify from LLMs was impossible, so I built a \textbf{production MCP server} with \textbf{75 tools and 10 resources}, enabling full Spotify control (playback, search, playlists, devices) from \textbf{Claude Desktop}.
    \item Setup and auth were a blocker, so I implemented \textbf{secure OAuth}, one-command Docker setup, and full developer docs for frictionless end-to-end usage.
    \item \textbf{Tech:} \textbf{MCP}, \textbf{Python 3.10+}, \textbf{Spotify Web API (OAuth2)}, \textbf{Docker}, \textbf{Claude Desktop}
\end{highlights}
\end{onecolentry}

        

    \section{Certifications}

        \begin{onecolentry}
            \begin{highlightsforbulletentries}
                \item \text{\textbf{\href{https://www.coursera.org/account/accomplishments/specialization/certificate/XTE2M8O9JST0}{Machine Learning Specialization}} - Stanford Online (Sep 2025).}
        
                \item \text{\textbf{Anthropic MCP Series} -} 
                \textbf{\href{http://verify.skilljar.com/c/fspc2qd4i7ks}{Introduction to MCPs}} \text{ and } 
                \textbf{\href{http://verify.skilljar.com/c/vjfcjaz3fjg7}{Advanced Topics in MCPs}} (Aug 2025).
        
                \item \text{\textbf{CS50 Series} - Harvard University}, 
                \textbf{\href{https://certificates.cs50.io/d2aeac53-b99f-41c4-a0b3-f62079869c87}{CS50AI}}, 
                \textbf{\href{https://cs50.harvard.edu/certificates/a52c49dc-ec4f-4810-9908-3bc189f64e54}{CS50X}}, 
                \textbf{\href{https://cs50.harvard.edu/certificates/2a78fd4e-2286-44c6-b04d-cd732e218372}{CS50P}} (Aug 2024 – Jun 2025).
            \end{highlightsforbulletentries}
        \end{onecolentry}
                  
    \section{Skills}

    \begin{onecolentry}
    \textbf{Languages:} Python (Advanced), C/C++, Java, JavaScript (ES6+), SQL, Bash, HTML5, CSS3.
    \end{onecolentry}
    
    \begin{onecolentry}
    \textbf{Tools \& Frameworks:} PyTorch, TensorFlow, Scikit-Learn, Hugging Face, LangGraph, Flask, FastAPI, React, Docker, AWS, Git.
    \end{onecolentry}
    
    \begin{onecolentry}
    \textbf{AI \& Platforms:} LLMs, RAG, AI Agents, MCP, REST APIs, OAuth, Webhooks.
    \end{onecolentry}

    


\end{document}




